\documentclass[10pt,a4paper]{amsart}
\usepackage[margin=19mm]{geometry}
\usepackage{amsmath,amssymb,mathtools}
\usepackage[scr=rsfs,cal=euler]{mathalfa}
\usepackage{xfrac}
\usepackage[unicode,hidelinks,bookmarks=false]{hyperref}
\newcommand{\ie}{i.e.\ }
\newcommand{\cf}{cf.\ }
\newcommand{\resp}{respectively\ }
\newcommand{\Subsection}{\subsection}
\providecommand{\nobreakdash}{\nobreak\hbox{-}}
\setlength{\emergencystretch}{2em}
\multlinegap0pt
\usepackage{amssymb}
\usepackage{mathtools}
\usepackage{enumerate}
\newtheorem{theorem}{Theorem}
\title{Inner approximations of doubling weights with applications to Beurling-Malliavin theory in Toeplitz kernels}
\author{Alex Bergman}
\date{Source: arXiv:2509.21229v4 (CC BY 4.0)}
\begin{document}
\maketitle

\noindent\emph{Source and licence.} Inner approximations of doubling weights with applications to Beurling-Malliavin theory in Toeplitz kernels. Version arXiv:2509.21229v4 (CC BY 4.0), distributed by arXiv under the Creative Commons Attribution 4.0 International licence, http://creativecommons.org/licenses/by/4.0/, as recorded in the arXiv licence metadata for this exact version and shown on the abstract page https://arxiv.org/abs/2509.21229v4. Full licence notice: \texttt{licenses/arxiv-2509.21229-CC-BY-4.0.txt}.\par\smallskip

\noindent\textit{Editorial scope.} Counted unit: the article's theorem thm:main (source0.tex lines 121--139). The article's complete original proof of it is delivered with it (source0.tex lines 602--607, one \texttt{proof} environment of 6 lines, which the article places away from the statement). No mathematics is written, repaired or re-derived: the statement and the proof are the article's own lines, in the article's own notation, and no retained line is substituted or re-flowed.  No further source-local context is needed: every cross-reference of every retained line resolves inside the two delivered spans.  Nothing was omitted from any retained span, so the package carries no omission record.  No retained line contains a citation, so the wrapper carries no bibliography and no bibliography entry.  No retained line contains a figure environment, an image-inclusion command or drawing code, so no image is delivered and the package declares no asset dependency.  Editorial formatting only, in addition: an AMS \texttt{amsart} preamble that restates the article's own declarations and macros used by the retained lines, namely the environments \texttt{theorem} on separate counters, together with the standard packages those lines need.  The retained environments print in this document as Theorem 1 in order of appearance, whereas the article prints the same items with the numbering of its own full document.  The complete native source of the article as distributed in the arXiv e-print for this exact version is bundled unchanged as \texttt{sources/185734/source0.tex}.
    \begin{theorem}\label{thm:main}
        Let $\alpha : \mathbb{R} \to \mathbb{R}$ be a smooth locally doubling weight satisfying
        \begin{equation*}
                \langle x \rangle^{-\kappa} \lesssim \alpha'(x), \text{ for some } \kappa < 1.
        \end{equation*}
        Then there exists a meromorphic inner function $J$ and a positive integer $N_{0}$, such that
        \begin{equation*}
            \lvert \alpha(x) -\arg J(x) \rvert \leq 2\pi, \text{ for all } x \in \mathbb{R},
        \end{equation*}
        and
        \begin{equation*}
            \lvert J'(x) \rvert \lesssim \alpha'(x)^{N_{0}}\langle x \rangle^{N_{0}}, \text{ for all }x \in \mathbb{R},
        \end{equation*}
        and
        \begin{equation*}
            \langle \lambda \rangle^{-N_{0}}\alpha'(\lambda)^{-N_{0}} \lesssim \lvert J'(\lambda) \rvert, \text{ for all } \lambda \in \Lambda,
        \end{equation*}
        where $\Lambda = \left\{ \lambda \in \mathbb{R} : \alpha(\lambda) = 0 \mod 2\pi\right\}$.
    \end{theorem}

    \begin{proof}[Proof of Theorem \ref{thm:main}]
            The result follows from the previous theorem after setting
            \begin{equation*}
                \Lambda = \left\{ \lambda \in \mathbb{R} : \alpha(\lambda) = 0 \mod 2\pi \right\}.
            \end{equation*}
    \end{proof}

\end{document}
